%% Correction audit de veracite - 2026-09-20 : citations de rapport de jury sans
%% support textuel dans le corpus (rapports 2018-2025) remplacees par des verbatims
%% verifies ou supprimees. Detail ci-dessous.
%% l.305 (ex) : la 1re phrase est conforme (RJ2020 l.2386) mais citee de memoire
%%   ("ne sait pas" au lieu de "ne sache pas") ; la suite ("Les resultats doivent etre
%%   encadres...") est un ajout INTROUVABLE -> verbatim RJ2020 l.2386 restitue, ajout retire.
\expandafter\ifx\csname CAPETINCLUDEMODE\endcsname\relax
\documentclass[11pt,a4paper]{article}
% preamble.tex — packages et commandes communes a toutes les fiches CAPET SII IE
\input{glyphtounicode}
\pdfgentounicode=1
\usepackage{cmap}
\usepackage[T1]{fontenc}
\usepackage[utf8]{inputenc}
\usepackage[french]{babel}
% babel[french] active a \begin{document} des raccourcis de ponctuation (:;!?)
% qui cassent la syntaxe to[...] de circuitikz (erreur "+ or - expected") ;
% on les desactive, ce qui ne change que l'espacement fin avant ces signes.
\AtBeginDocument{\shorthandoff{:;!?}}
\usepackage[a4paper, top=1.6cm, bottom=1.6cm, left=1.4cm, right=1.4cm]{geometry}
\setlength{\headheight}{14.5pt}
\usepackage{amsmath, amssymb, amsfonts}
\usepackage{siunitx}
\sisetup{locale=FR, output-decimal-marker={,}, inter-unit-product=\ensuremath{{}\cdot{}},
         group-minimum-digits=4}
% Unites non fournies par siunitx, declarees ici une fois pour toutes (fiches + exercices).
% Les gardes \ifdefined evitent tout conflit avec une redeclaration locale restee dans un
% fichier : ne pas retirer les gardes, ne pas redeclarer ces unites dans les fichiers.
\DeclareSIUnit{\jour}{j}
\ifdefined\var\else\DeclareSIUnit{\var}{var}\fi   % puissance reactive
\ifdefined\vad\else\DeclareSIUnit{\vad}{VAD}\fi   % puissance deformante
\ifdefined\tr\else\DeclareSIUnit{\tr}{tr}\fi      % tour (tr/min)
\usepackage{array, booktabs, multirow, tabularx}
\usepackage{xcolor}
\usepackage[most, breakable]{tcolorbox}
\usepackage{multicol}
\usepackage{enumitem}
\usepackage{mathtools}
\usepackage{lmodern}
% Refonte visuelle 2026 (direction "moderne coloré") : Lato pour le texte,
% maths sans-serif assorties via newtxsf (charge apres amsmath/amssymb).
\usepackage[default]{lato}
\usepackage{newtxsf}
\usepackage{microtype}
\usepackage{fancyhdr}
\usepackage{tikz}
\usetikzlibrary{positioning,arrows.meta}
\usepackage{circuitikz}
\usepackage{lastpage}

% ── Couleurs par domaine (palette modernisée 2026) ────────────────────
\definecolor{colElec}{RGB}{29, 78, 216}        % Électrotechnique — bleu vif
\definecolor{colElecLight}{RGB}{219, 234, 254} % fond bleu clair
\definecolor{colAuto}{RGB}{21, 128, 61}        % Automatique — vert
\definecolor{colMeca}{RGB}{194, 65, 12}        % Mécanique — orange
\definecolor{colTherm}{RGB}{185, 28, 28}       % Thermique — rouge
\definecolor{colQT}{RGB}{232, 244, 233}        % (hérité, plus utilisé par questiontype)
\definecolor{colJury}{RGB}{255, 248, 225}      % fond jury (crème)
\definecolor{colJuryFrame}{RGB}{202, 138, 4}   % bordure/bandeau jury (ambre)
\definecolor{colGrayLight}{RGB}{245, 245, 247} % fond gris neutre
\definecolor{colSeparator}{RGB}{180, 180, 180} % gris séparateur
\definecolor{colExo}{RGB}{8, 145, 178}          % Banque d'exercices — cyan (distinct des fiches)
% ── Encre et gris de la refonte ───────────────────────────────────────
\definecolor{colInk}{RGB}{23, 32, 42}          % texte courant (quasi-noir)
\definecolor{colSlate}{RGB}{83, 98, 118}       % gris ardoise (méta, en-têtes)
\definecolor{colRule}{RGB}{203, 213, 225}      % filets clairs
\definecolor{colJuryBg}{RGB}{255, 248, 225}    % fond crème jury
\definecolor{colPresence}{RGB}{83, 98, 118}    % bloc 5 (présence dans les sujets)
\AtBeginDocument{\color{colInk}}

% ── Style encadré exercice + corrigé (banque d'exercices) ─────────────
\tcbset{
  exoBox/.style={
    breakable,
    enhanced,
    colback=colExo!8!white,
    colframe=colExo,
    colbacktitle=colExo,
    coltitle=white,
    fonttitle=\bfseries\small,
    left=4pt, right=4pt, top=3pt, bottom=3pt,
    before skip=5pt, after skip=5pt,
  }
}

% ── Bascule enonce / corrige (banque d'exercices) ─────────────────────
% \CAPETCORRIGE est defini en ligne de commande (pdflatex "\def\CAPETCORRIGE{1}\input{...}")
% pour compiler la variante corriges d'un meme fichier source ; par defaut (non defini),
% la compilation produit la variante enonces. \CAPETSHOW{enonce}{corrige} choisit la branche.
%
% TROIS MODES (refonte 2026-09-20, manuel de revision) :
%   - aucun drapeau   -> imprime l'enonce            (variante enonces, INCHANGE)
%   - \CAPETCORRIGE   -> imprime le corrige          (variante corriges, INCHANGE)
%   - \CAPETMANUEL    -> imprime l'enonce ET met le corrige de cote pour l'annexe
%
% Le mode manuel stocke chaque corrige dans une macro (\csgdef), sans ecriture
% fichier ni re-tokenisation : le contenu (circuitikz, accents, %) traverse intact.
% L'annexe les rejoue ensuite dans l'ordre via \CAPETANNEXE.
\usepackage{etoolbox}

\newcounter{CAPETexo}          % numero continu d'exercice dans le manuel
\setcounter{CAPETexo}{0}
\newcounter{CAPETannexe}       % compteur de relecture en annexe

% Titre courant de la notion en cours (pose par \CAPETNOTION), sert d'intertitre annexe
\newcommand{\CAPETnotioncode}{}
\newcommand{\CAPETnotiontitre}{}

% \CAPETNOTION{code}{titre} : ouvre une notion dans le manuel (memorise le libelle,
% et marque le rang du premier exercice de la notion pour l'intertitre d'annexe).
\newcommand{\CAPETNOTION}[2]{%
  \renewcommand{\CAPETnotioncode}{#1}%
  \renewcommand{\CAPETnotiontitre}{#2}%
  \csgdef{CAPETnotionstart\the\numexpr\value{CAPETexo}+1\relax}{#1 -- #2}%
}

% ── \degree robuste ───────────────────────────────────────────────────
% Certains fichiers d'exercices utilisent \degree. Ni le preambule ni siunitx
% ne le definissent, et \providecommand local ne suffit pas : sous hyperref,
% \degree existe deja comme commande factice (\ERROR), donc \providecommand
% s'abstient et l'appel echoue a la compilation du manuel (erreur invisible
% en compilation isolee). On le (re)definit donc inconditionnellement ici,
% apres le chargement des packages, pour tout le corpus.
\AtBeginDocument{\providecommand{\degree}{}\renewcommand{\degree}{\ensuremath{^\circ}}}

% Libelle decoratif : les fichiers d'exercices utilisent aussi \CAPETSHOW dans
% leur \fancyhead et leur bloc de titre, pour afficher "(enonces)" ou "(corriges)".
% Ces appels-la ne sont PAS des exercices : en mode manuel ils ne doivent ni
% consommer un numero, ni alimenter l'annexe. On les reconnait a leur premier
% argument, strictement egal au libelle ci-dessous.
\newcommand{\CAPETLIBELLEENONCES}{(\'enonc\'es)}

\ifdefined\CAPETMANUEL
  % --- Mode manuel : enonce en place, corrige differe vers l'annexe ---
  \newcommand{\CAPETSHOW}[2]{%
    \def\CAPETargun{#1}%
    \ifx\CAPETargun\CAPETLIBELLEENONCES
      #1%            appel decoratif : rendu tel quel, hors numerotation
    \else
      \CAPETSHOWEXO{#1}{#2}%
    \fi
  }
  \newcommand{\CAPETSHOWEXO}[2]{%
    \stepcounter{CAPETexo}%
    \csgdef{CAPETcorrige\theCAPETexo}{#2}%
    \edef\CAPETtmpnum{\theCAPETexo}%
    #1%
    \par\vspace{2pt}%
    {\footnotesize\itshape\color{colExo}%
     \hyperref[capet:cor:\CAPETtmpnum]{Corrige de l'exercice \CAPETtmpnum\ $\rightarrow$ p.~\pageref{capet:cor:\CAPETtmpnum}}%
     \label{capet:eno:\CAPETtmpnum}}%
  }
  % \CAPETANNEXE : rejoue tous les corriges captures, dans l'ordre, avec intertitres
  \newcommand{\CAPETANNEXE}{%
    \setcounter{CAPETannexe}{0}%
    \loop\ifnum\value{CAPETannexe}<\value{CAPETexo}%
      \stepcounter{CAPETannexe}%
      \edef\CAPETcurnum{\theCAPETannexe}%
      \ifcsname CAPETnotionstart\CAPETcurnum\endcsname
        \par\medskip
        {\normalsize\bfseries\color{colExo}\csuse{CAPETnotionstart\CAPETcurnum}}%
        \par\nobreak\smallskip
      \fi
      \begin{tcolorbox}[exoBox, breakable,
        title={Corrige de l'exercice \CAPETcurnum\
               \hfill\normalfont\footnotesize
               \hyperref[capet:eno:\CAPETcurnum]{enonce p.~\pageref{capet:eno:\CAPETcurnum}}}]%
      \label{capet:cor:\CAPETcurnum}%
      \csuse{CAPETcorrige\CAPETcurnum}%
      \end{tcolorbox}%
    \repeat
  }
\else
  % --- Modes historiques, comportement strictement inchange ---
  \newcommand{\CAPETSHOW}[2]{\ifdefined\CAPETCORRIGE #2\else #1\fi}
  \newcommand{\CAPETANNEXE}{}
\fi

% ── Difficulté 1 à 3 (étoiles pleines, amssymb) ───────────────────────
\newcommand{\difficulte}[1]{%
  \ifcase#1\relax
  \or $\bigstar$\;niveau 1/3%
  \or $\bigstar\bigstar$\;niveau 2/3%
  \or $\bigstar\bigstar\bigstar$\;niveau 3/3%
  \fi
}

% ── Style encadré extrait de document technique (banque d'exercices) ──
\tcbset{
  docTech/.style={
    breakable,
    enhanced,
    colback=colGrayLight,
    colframe=colSeparator,
    fonttitle=\bfseries\small\itshape,
    coltitle=black,
    colbacktitle=colGrayLight,
    left=4pt, right=4pt, top=3pt, bottom=3pt,
    before skip=4pt, after skip=4pt,
  }
}

% ── Style tcolorbox de section principale ────────────────────────────
% Refonte 2026 : bandeau titre coloré arrondi + fond très pâle, sans cadre.
\tcbset{
  ficheSection/.style={
    breakable, enhanced jigsaw,
    frame hidden, boxrule=0pt, arc=1.8mm,
    colback=#1!4!white,
    colbacktitle=#1,
    coltitle=white,
    fonttitle=\sffamily\bfseries,
    left=7pt, right=7pt, top=5pt, bottom=5pt,
    toptitle=3pt, bottomtitle=3pt,
    before skip=10pt, after skip=8pt,
  }
}

% ── Style encadré question type ──────────────────────────────────────
% Carte blanche, bandeau sombre, liseré très clair.
\tcbset{
  questiontype/.style={
    breakable, enhanced jigsaw,
    colback=white,
    frame hidden, boxrule=0pt, arc=1.8mm,
    colbacktitle=colElec!78!black,
    coltitle=white,
    fonttitle=\sffamily\bfseries\small,
    left=7pt, right=7pt, top=4pt, bottom=4pt,
    toptitle=2.5pt, bottomtitle=2.5pt,
    before skip=8pt, after skip=8pt,
    borderline={0.6pt}{0pt}{colElec!12!white},
  }
}

% ── Style encadré jury ───────────────────────────────────────────────
% Bandeau ambre, fond crème, sans cadre.
\tcbset{
  jury/.style={
    breakable, enhanced jigsaw,
    colback=colJuryBg,
    frame hidden, boxrule=0pt, arc=1.8mm,
    colbacktitle=colJuryFrame,
    coltitle=white,
    fonttitle=\sffamily\bfseries,
    left=7pt, right=7pt, top=4pt, bottom=4pt,
    toptitle=3pt, bottomtitle=3pt,
    before skip=10pt, after skip=8pt,
  }
}

% ── En-tête / pied de page communs (refonte 2026) ─────────────────────
% #1 code fiche, #2 titre court, #3 couleur domaine
\newcommand{\ficheHeaders}[3]{%
  \fancyhead[L]{\small\sffamily{\bfseries\textcolor{#3}{#1}}\hspace{0.7em}\textcolor{colSlate}{#2}}%
  \fancyhead[R]{\small\sffamily\textcolor{colSlate}{CAPET SII IE\, --- \,Session 2026}}%
  \fancyfoot[C]{\small\sffamily\textcolor{colSlate}{\thepage\,/\,\pageref{LastPage}}}%
  \fancyhead[C]{}\fancyfoot[L]{}\fancyfoot[R]{}%
}

% ── Bloc en-tête de fiche (BLOC 0, refonte 2026) ──────────────────────
% #1 code  #2 titre  #3 couleur domaine  #4 domaine  #5 ligne présence  #6 sous-thèmes
\newcommand{\ficheEntete}[6]{%
  \begin{tcolorbox}[enhanced, breakable,
    frame hidden, boxrule=0pt, arc=2.5mm,
    interior style={left color=#3, right color=#3!75!black},
    left=9pt, right=9pt, top=7pt, bottom=7pt,
    before skip=0pt, after skip=6pt]
    {\color{white}\sffamily\bfseries\LARGE #1\quad #2}\\[4pt]
    {\color{#3!22!white}\sffamily\small \MakeUppercase{#4}\quad\textbar\quad #5}\\[1pt]
    {\color{#3!35!white}\sffamily\footnotesize Sous-thèmes : #6}
  \end{tcolorbox}%
}

% ── Mise en page ─────────────────────────────────────────────────────
\pagestyle{fancy}
\fancyhf{}
\renewcommand{\headrulewidth}{0.5pt}
\renewcommand{\footrulewidth}{0pt}
\AtBeginDocument{%
  \renewcommand{\headrule}{{\color{colRule}\hrule width\headwidth height 0.5pt}}%
}

\setlist[itemize]{leftmargin=1.2em, itemsep=1pt, topsep=2pt, parsep=0pt}
\setlist[enumerate]{leftmargin=1.2em, itemsep=3pt, topsep=2pt, parsep=0pt}

\setlength{\parskip}{2pt}
\setlength{\columnsep}{10pt}

\def\CAPETBEGINFILE{\begin{document}}
\def\CAPETENDFILE{\end{document}}
\else
\def\CAPETBEGINFILE{}
\def\CAPETENDFILE{}
\fi

\ficheHeaders{T22}{CAN/CNA - Numérisation}{colElec}

\CAPETBEGINFILE

% -----------------------------------------------------------------------
%  BLOC EN-TETE
% -----------------------------------------------------------------------

\ficheEntete{T22}{CAN/CNA - Numérisation et échantillonnage}{colElec}%
  {électronique - Chaîne d'acquisition numérique}%
  {\textbf{Présence :} \textbf{6/13} sujets (2013-2026) - thème \emph{fréquent}}%
  {quantum - nombre de bits - codage binaire/hex - théorème de Shannon - filtre anti-repliement - chaîne capteur→conditionneur→CAN - CNA en sortie de correcteur}

% -----------------------------------------------------------------------
%  1. RAPPELS DE COURS
% -----------------------------------------------------------------------

\begin{tcolorbox}[
  ficheSection=colElec,
  title={1. Rappels de cours},
]

% -- 1.1 CAN - Quantum et resolution -----------------------------------
\textbf{1.1 Convertisseur Analogique-Numérique (CAN) - Quantum}

\[
  q = \frac{V_{ref}}{2^n - 1} \approx \frac{V_{ref}}{2^n} \quad (n \geq 8)
  \qquad
  N = \frac{V_E}{q} = \frac{V_E \cdot 2^n}{V_{ref}}
  \qquad
  \text{Erreur max :} \pm\frac{q}{2}
\]

\smallskip
\textbf{Numérisation par pas de quantum (courbe analogique et escalier numérique) :}
\begin{center}
\begin{tikzpicture}[scale=0.85, font=\footnotesize]
  \draw[->] (0,0) -- (6.3,0) node[right] {$t$};
  \draw[->] (0,-0.3) -- (0,3.2) node[above] {$V$};
  \draw[thick, colElec!45] (0,0.3) -- (6,2.9);
  \draw[thick, colElec] (0,0.3) -- (1,0.3) -- (1,0.85) -- (2,0.85) -- (2,1.4) -- (3,1.4) -- (3,1.95) -- (4,1.95) -- (4,2.5) -- (5,2.5) -- (5,2.85) -- (6,2.85);
  \draw[<->] (6.2,1.95) -- (6.2,2.5);
  \node[right] at (6.25,2.2) {\scriptsize $q$};
\end{tikzpicture}
\end{center}
Le signal analogique (clair) est approché par des paliers de largeur $q$ (foncé) : chaque palier correspond à un code $N$, avec une erreur de quantification bornée à $\pm q/2$.

\textbf{Nombre de bits nécessaires} pour une précision $\delta$ sur la plage $V_{ref}$ :
\[
  2^n \geq \frac{V_{ref}}{\delta}
  \qquad \Rightarrow \qquad
  n \geq \log_2\!\left(\frac{V_{ref}}{\delta}\right)
\]

\begin{center}
\renewcommand{\arraystretch}{1.25}
\begin{tabular}{cccc}
\toprule
$n$ (bits) & $2^n$ & $q$ pour $V_{ref}=5\,\mathrm{V}$ & $q$ pour $V_{ref}=10\,\mathrm{V}$ \\
\midrule
8  & 256   & $\approx 19{,}6\,\mathrm{mV}$ & $\approx 39{,}2\,\mathrm{mV}$ \\
10 & 1\,024 & $\approx 4{,}9\,\mathrm{mV}$  & $\approx 9{,}8\,\mathrm{mV}$ \\
12 & 4\,096 & $\approx 1{,}2\,\mathrm{mV}$  & $\approx 2{,}44\,\mathrm{mV}$ \\
\bottomrule
\end{tabular}
\end{center}

\medskip\hrule\vspace{4pt}

% -- 1.2 CNA ----------------------------------------------------------
\textbf{1.2 Convertisseur Numérique-Analogique (CNA)}

\[
  V_S = N \cdot q = N \cdot \frac{V_{ref}}{2^n}
\]
Utilisé en sortie du correcteur numérique dans une boucle d'asservissement échantillonnée.

\medskip\hrule\vspace{4pt}

% -- 1.3 Shannon ------------------------------------------------------
\textbf{1.3 Théorème de Shannon - échantillonnage}

\[
  f_e \geq 2\,f_{max}
\]
\begin{itemize}
  \item Filtre anti-repliement (passe-bas) avant le CAN : $f_c \leq f_e/2$
  \item Période d'échantillonnage en asservissement : $T_e \leq T_{systeme}/10$ (règle pratique)
  \item Aliasing (repliement spectral) si $f_e < 2\,f_{max}$ - signal reconstruit corrompu
\end{itemize}

\medskip\hrule\vspace{4pt}

% -- 1.4 Chaene d'acquisition -----------------------------------------
\textbf{1.4 Chaîne d'acquisition complète}

\[
  \text{Grandeur physique} \xrightarrow{\text{Capteur}} v_T
  \xrightarrow{\text{Conditionneur+AMP}} v_A \in [0\,;\,V_{ref}]
  \xrightarrow{\text{CAN}} G_N
\]
La tension $v_A$ doit couvrir \textbf{toute} la plage $[0\,;\,V_{ref}]$ pour utiliser au maximum la résolution.

\textbf{Sensibilité globale de la chaîne :}
\[
  S = \frac{\Delta G_N}{\Delta x} \quad [\text{LSB/unité physique}]
\]

\medskip\hrule\vspace{4pt}

% -- 1.5 Codage hexadecimal --------------------------------------------
\textbf{1.5 Codage hexadécimal}

Grouper par 4 bits depuis le bit de poids faible (LSB) : ex. $1842_{10} = 0111\,0011\,0010_2 = 732_{16}$.

\medskip\hrule\vspace{4pt}

% Complément du référentiel - 2026-11-10 (sans TikZ : les ids de schémas restent stables)
\textbf{1.6 Structures de CAN et de CNA (complément)}

\smallskip
\textbf{CAN flash} : $2^n - 1$ comparateurs en parallèle, conversion en un coup d'horloge, très rapide mais limité à 8 à 10 bits. \textbf{Approximations successives} (SAR) : un CNA interne et un comparateur fixent les bits un par un, du poids fort au poids faible, en $t_{conv} \approx n\,T_{clk}$ ; c'est le compromis courant (10 à 16 bits). \textbf{Double rampe} : intégration de $V_e$ pendant $T_1$ fixe, puis de $-V_{ref}$ jusqu'au retour à zéro en $T_2$, d'où $V_e = V_{ref}\,T_2/T_1$ ; lent, précis, il rejette le 50\,Hz (multimètres). \textbf{Sigma-delta} : sur-échantillonnage et filtrage numérique, 16 à 24 bits à débit modéré. \textbf{CNA à réseau R-2R} : $|V_s| = V_{ref}\,N/2^n$ avec deux valeurs de résistances seulement.

\end{tcolorbox}

% -----------------------------------------------------------------------
%  2. FORMULAIRE SYNTHETIQUE
% -----------------------------------------------------------------------

\begin{tcolorbox}[
  ficheSection=colElec,
  title={2. Formulaire synthétique},
]
\small
\begin{multicols}{2}
\renewcommand{\arraystretch}{1.45}
\begin{tabular}{@{}p{4.8cm}ll@{}}
\toprule
\textbf{Formule} & \textbf{Usage} & \textbf{U.} \\
\midrule
\multicolumn{3}{@{}l}{\footnotesize\textit{CAN}} \\
$q = V_{ref}/(2^n-1)$ & Quantum exact & V \\
$q \approx V_{ref}/2^n$ & Quantum approx. & V \\
$N = V_E/q$ & Code numérique & LSB \\
$n \geq \log_2(V_{ref}/\delta)$ & Nbre de bits & - \\
Erreur max $= \pm q/2$ & Quantification & V \\
\midrule
\multicolumn{3}{@{}l}{\footnotesize\textit{CNA}} \\
$V_S = N \cdot V_{ref}/2^n$ & Reconstruction & V \\
\bottomrule
\end{tabular}

\columnbreak

\renewcommand{\arraystretch}{1.45}
\begin{tabular}{@{}p{4.8cm}ll@{}}
\toprule
\textbf{Formule} & \textbf{Usage} & \textbf{U.} \\
\midrule
\multicolumn{3}{@{}l}{\footnotesize\textit{Shannon}} \\
$f_e \geq 2\,f_{max}$ & Échantillonnage & Hz \\
$f_c \leq f_e/2$ & Filtre AR & Hz \\
$T_e \leq T_{syst}/10$ & Asservissement & s \\
\midrule
\multicolumn{3}{@{}l}{\footnotesize\textit{Données types}} \\
10 bits / 5\,V : $q \approx 4{,}9\,\mathrm{mV}$ & PIC 18F & - \\
12 bits / 10\,V : $q \approx 2{,}44\,\mathrm{mV}$ & Précision & - \\
S7-1500 : PIW8 $\in [0\,;\,27\,648]$ pour 4-20\,mA & API & - \\
\bottomrule
\end{tabular}
\end{multicols}
\normalsize
\end{tcolorbox}

% -----------------------------------------------------------------------
%  3. QUESTIONS TYPES
% -----------------------------------------------------------------------

\begin{tcolorbox}[
  ficheSection=colElec,
  title={3. Questions types - extraites des sujets réels},
]

% -- QT1 -------------------------------------------------------------

\begin{tcolorbox}[questiontype,
  title={QT1 - Calcul du quantum et codage hexadécimal \hfill\normalfont\footnotesize E1 2013 (tour laboratoire)}
]
\textbf{Énoncé.}
Signal de température compris entre $-50^\circ$C et $+50^\circ$C, précision $0{,}1^\circ$C, tension de sortie du capteur-conditionneur : $V \in [0\,;\,10]\,\mathrm{V}$. Déterminer le nombre de bits du CAN et calculer le mot hexadécimal pour $\theta = -5^\circ$C.

\tcbline

\textbf{Corrigé.}

\textbf{Étape 1 — Précision en tension.} $\delta = \dfrac{0{,}1}{100} \times 10 = 10\,\mathrm{mV}$.

\textbf{Étape 2 — Nombre de bits.}
\[
  2^n \geq \frac{10}{0{,}01} = 1000 \quad \Rightarrow \quad n \geq 10 \; \Rightarrow \; \text{choix : } \mathbf{n = 12\,\text{bits}}
\]
Quantum : $q = 10/(4095) \approx 2{,}44\,\mathrm{mV}$.

\textbf{Étape 3 — Mot numérique pour $\theta = -5^\circ$C.} $V_E = \dfrac{10}{100}(-5+50) = 4{,}5\,\mathrm{V}$.
\[
  N = \frac{4{,}5}{2{,}44 \times 10^{-3}} = 1842_{10} \quad = \quad 0111\,0011\,0010_2 \quad = \quad \mathbf{732_{16}}
\]

\textbf{Conclusion.} Le choix $n = 12$ bits offre une marge par rapport au minimum théorique de 10 bits : cette marge absorbe les non-linéarités du CAN réel et garantit la précision de $0{,}1^\circ$C exigée sur toute la plage.
\end{tcolorbox}

\vspace{4pt}

% -- QT2 -------------------------------------------------------------

\begin{tcolorbox}[questiontype,
  title={QT2 - Sensibilité de la chaîne et mot numérique \hfill\normalfont\footnotesize E2 BAC3 2026 (Q16-Q17)}
]
\textbf{Énoncé.}
Capteur Jumo + amplificateur/conditionneur + CAN 8 bits (ADS7029-Q1), pleine échelle $\SI{3}{\volt}$, plage $T \in [-5\,;\,90^\circ\mathrm{C}]$, avec $v_A = c \cdot T$ et $c = \SI{0,0316}{\volt\per\degreeCelsius}$. Montrer que la sensibilité globale vaut $S = 2{,}684\,\mathrm{LSB/^\circ C}$, puis en déduire le mot numérique $N_C$ pour $T_C = 55^\circ$C.

\tcbline

\textbf{Corrigé.} \textit{[corrigé reconstitué]}

\textbf{Étape 1 — Quantum du CAN 8 bits.} $q = \dfrac{3}{2^8 - 1} = \dfrac{3}{255} \approx \SI{11,76}{\milli\volt}$.

\textbf{Étape 2 — Vérification de la sensibilité.} Plage de $95^\circ$C couverte par $\SI{3}{\volt}$ : $c = 3/95 \approx \SI{0,0316}{\volt\per\degreeCelsius}$ ; d'où
\[
  S = \frac{c}{q} = \frac{0{,}0316}{11{,}76 \times 10^{-3}} = \frac{255}{95} \approx 2{,}684\,\mathrm{LSB/^\circ C}
\]

\textbf{Étape 3 — Mot numérique.} Comme $v_A = c \cdot T$, pour $T_C = 55^\circ$C :
\[
  N_C = S \times T_C = 2{,}684 \times 55 \approx 147{,}6 \quad \Rightarrow \quad N_C \approx \mathbf{148\,\mathrm{LSB}}
\]
En hexadécimal : $148_{10} = 1001\,0100_2 = \mathrm{94_{16}}$.

\textbf{Conclusion.} $N_C \approx 148\,\mathrm{LSB}$ reste dans la plage codable sur 8 bits (0 à 255), et même à $T_{max} = 90^\circ$C le code $\approx 242$ ne sature pas : la chaîne d'acquisition exploite au mieux la pleine échelle sans risque de saturation.
\end{tcolorbox}

\vspace{4pt}

% -- QT3 -------------------------------------------------------------

\begin{tcolorbox}[questiontype,
  title={QT3 - Théorème de Shannon - filtre anti-repliement \hfill\normalfont\footnotesize E1 2026 (B17)}
]
\textbf{Énoncé.}
Rappeler le théorème de Shannon. Définir le rôle du filtre anti-repliement. Déduire sa fréquence de coupure $f_c$ si la fréquence d'échantillonnage est $f_e$.

\tcbline

\textbf{Corrigé.}

\textbf{Étape 1 — Théorème de Shannon.} Pour reconstituer sans aliasing un signal de fréquence maximale $f_{max}$, la fréquence d'échantillonnage $f_e$ doit vérifier $f_e \geq 2\,f_{max}$.

\textbf{Étape 2 — Rôle du filtre anti-repliement.} Filtre passe-bas placé \textbf{avant le CAN}, qui élimine toutes les composantes spectrales supérieures à $f_e/2$. Sans ce filtre, ces composantes se replient dans la bande utile (aliasing) et corrompent le signal numérique.

\textbf{Étape 3 — Fréquence de coupure (filtre parfait).}
\[
  f_c = \frac{f_e}{2}
\]
En pratique, $f_c < f_e/2$ pour tenir compte de la transition finie du filtre réel.

\textbf{Conclusion.} Sans ce filtre anti-repliement correctement réglé, le respect théorique de $f_e \geq 2f_{max}$ ne suffit pas à garantir un signal numérisé fidèle : c'est la combinaison des deux qui satisfait le cahier des charges d'acquisition.
\end{tcolorbox}

\end{tcolorbox}

% -----------------------------------------------------------------------
%  4. ATTENTES DU JURY
% -----------------------------------------------------------------------

\begin{tcolorbox}[
  jury,
  title={4. Attentes du jury - Erreurs récurrentes et points de vigilance},
]

\begin{enumerate}[label=\textbf{\arabic*.}, leftmargin=*, itemsep=5pt]

  \item \textbf{Maîtriser la formule du quantum} et la distinguer de la précision physique. Résolution $= q$ en volt ; précision physique $= q / (\text{sensibilité capteur})$ en unité mesurée.

  \item \textbf{Utiliser $2^n - 1$ ou $2^n$} selon le contexte. Pour $n \geq 8$, l'approximation $2^n$ est admise mais doit être justifiée.

  \item \textbf{Shannon : ne pas se limiter à l'énoncé.} Appliquer la conséquence immédiate : $f_c \leq f_e/2$ pour le filtre anti-repliement, $T_e \leq T_{syst}/10$ pour l'asservissement.

  \item \textbf{Filtre anti-repliement obligatoire.} Le jury attend sa mention dès qu'un CAN apparaît dans une boucle de mesure - même si la question ne le demande pas explicitement.

  \item \textbf{Codage hexadécimal.} Grouper 4 bits depuis le LSB, pas depuis le MSB. Vérifier la cohérence : $N_{10} \times q$ doit retrouver $V_E$.

  \item \textbf{Confondre CAN et CNA.} Le CAN est en voie de mesure (entrée analogique → code numérique) ; le CNA est en voie de commande (sortie correcteur numérique → actionneurs analogiques).

\end{enumerate}

\medskip
\begin{tcolorbox}[
  colback=colJuryFrame!12!colJury,
  colframe=colJuryFrame,
  left=4pt, right=4pt, top=2pt, bottom=2pt,
  fontupper=\small\itshape
]
Rapport jury 2020 : «~Le jury déplore que la majeure partie des candidats ne sache pas déterminer la résolution d'un capteur.~»
\end{tcolorbox}

\end{tcolorbox}

% -----------------------------------------------------------------------
%  PIED - PRESENCE DANS LES SUJETS
% -----------------------------------------------------------------------

\begin{tcolorbox}[
  ficheSection=colPresence,
  title={Présence dans les sujets analysés (2013-2026)},
]
\begin{center}
\footnotesize\renewcommand{\arraystretch}{1.15}
\begin{tabular}{lccccccccccccc}
\toprule
\textbf{Année} & 13 & 14 & 15 & 16 & 17 & 18 & 19 & 20 & 21 & 22 & 23 & 24 & 25 \\
\midrule
E1 (commune)   & $\checkmark$ & - & - & - & - & - & - & - & - & $\checkmark$ & - & - & $\checkmark$ \\
E2 discipl. IE & - & - & - & - & - & $\checkmark$ & - & $\checkmark$ & $\checkmark$ & $\checkmark$ & - & - & $\checkmark$ \\
E2 appl. comm. & - & - & - & - & - & - & - & - & - & - & - & - & - \\
\bottomrule
\end{tabular}
\end{center}
\smallskip
\footnotesize\textbf{Sous-thèmes fréquents :}
calcul du quantum - choix du nombre de bits - codage hexa - théorème de Shannon - filtre anti-repliement - chaîne capteur→CAN complète - sensibilité globale

\medskip
\textbf{Ordres de grandeur (repères rapides).}
\begin{itemize}[leftmargin=1.3em]
  \item Réseau BT FR : tension phase-neutre typique \SI{230}{\volt}, tension phase-phase typique \SI{400}{\volt}, fréquence \SI{50}{\hertz}.
  \item Courants d'étude fréquents en sujet : de quelques \si{\ampere} à quelques dizaines de \si{\ampere}.
  \item En conversion d'énergie, un rendement attendu reste généralement entre 0{,}8 et 0{,}98 selon la technologie.
\end{itemize}

\medskip
\textbf{Synthèse express (à refaire sans support).}
\begin{enumerate}[leftmargin=1.4em]
  \item Poser les hypothèses du modèle et les conventions de signe.
  \item Écrire 2 à 4 relations clés avec unités.
  \item Réaliser un calcul numérique complet et vérifier l'ordre de grandeur.
  \item Justifier le résultat par un argument technique de niveau CAPET.
\end{enumerate}

\end{tcolorbox}

\CAPETENDFILE
